% ---- pg-536 （书页 524） ----

\chapter{聚类、距离方法与多维标度变换}
\phantomsection\label{chap:12}

\section{引言}
\phantomsection\label{sec:12.1}

初等的、探索性的方法对理解多元关系的复杂性质往往很有用. 例如, 我们在本书中自始至终一直强调数据图的价值. 在本章中, 我们将补充一些建立在距离度量基础上的数据表示方法, 这种方法有点像数据的逐步分组法则(算法). 在数据中寻找某种“自然的”分组结构是一项重要的探索性技术. 分组可以提供一些非正式的方法, 以帮助我们估计维数, 确定离群值, 以及提出涉及各种关系的假设.

分组或聚类与上一章中讨论的分类有区别. 分类问题中组的数目是已知的, 其操作目标是将一个新的对象分派给这些组之一. 聚类分析是一种比较原始的技术, 对组的数目或组的结构不做任何假定. 分类是在相似性或距离(非相似性)的基础上进行的. 所需的输入是一些相似性度量, 或者说是一些可从中算出相似性的数据.

为了说明在定义一种自然分组方法时所遇到的困难的实质, 我们来考虑如何将一副普通纸牌中四种花色的 A, K, Q, J 共 16 张牌进行分类. 我们将某些分组情况示于\hyperref[fig:12.1]{图 12.1}. 见到这些图之后立刻就会明白: 有意义的分组依赖于对“相似”一词所下的定义.

\begin{mvfigure}
\mvinc{pg-536}{0.252,0.594,0.789,0.764}
\phantomsection\label{fig:12.1}\mvcap{图 12.1\quad 纸牌分组}
\end{mvfigure}

在聚类分析的大多数实际应用中, 研究者清楚了解如何区分“好”的分组与“坏”的分组. 那么为什么不列出有可能的分组, 然后从中选择“最优”分组, 以供未来研究?
